% Generated draft: requires mathematical and source review.
\documentclass[11pt]{article}
\usepackage[margin=25mm]{geometry}
\usepackage{fontspec}
\setmainfont{Latin Modern Roman}
\usepackage{amsmath,amssymb,mathtools,unicode-math}
\setmathfont{Latin Modern Math}
\usepackage{enumitem,needspace,xurl,hyperref}
\setlength{\parindent}{0pt}
\setlength{\parskip}{5pt}
\setlength{\emergencystretch}{3em}
\newcommand{\N}{\mathbb{N}}
\newcommand{\Z}{\mathbb{Z}}
\newcommand{\Q}{\mathbb{Q}}
\newcommand{\R}{\mathbb{R}}
\newcommand{\C}{\mathbb{C}}
\newcommand{\F}{\mathbb{F}}
\newcommand{\A}{\mathbb{A}}
\newcommand{\PP}{\mathbb{P}}
\newcommand{\E}{\mathbb{E}}
\newcommand{\eps}{\varepsilon}
\newcommand{\dd}{\,\mathrm{d}}
\begin{document}
\section*{UnsolvedMath ID: 20003301}
\textbf{Source problem code:} catalogue entry
\subsection{Definitions and mathematical statement}
For $n\ge1$ and $v>0$, let
\[
 W_v=\left\{y\in\R^n:|q\cdot y-p|<\|q\|_\infty^{-v}
 \text{ for infinitely many }(q,p)\in(\Z^n\setminus\{0\})\times\Z\right\}.
\]
This is dual approximation by integer linear forms, using the max norm, as defined in companion permanent ID $20003299$. A vector is very well approximable if it belongs to $W_{n+\varepsilon}$ for some $\varepsilon>0$. For a smooth $d$-dimensional nondegenerate submanifold $\mathcal M\subset\R^n$, D. Kleinbock's Question 26 asks for estimates of
\[
 \dim_H(\mathcal M\cap W_v)
\]
using dynamics. In the source's chart formulation, nondegeneracy means that positive-order partial derivatives of a local parametrization span $\R^n$ at the points under consideration. The exact differentiability/order and exceptional-point conventions should be checked in the source references. The question covers $v>0$; the nontrivial very-well-approximable regime is $v>n$.

The attached aid reports that the expected value
\[
 d-1+\frac{n+1}{v+1}\qquad(v>n)
\]
has been established in substantial classes, including the Veronese-curve value $(n+1)/(v+1)$, but not for every nondegenerate manifold. It supplies no complete primary citations or class hypotheses, so no general exact-dimension theorem is asserted here.

Its separate elementary model is: for every compact interval $I\subset\R$ with nonempty interior and every $v>0$,
\[
 \dim_H\{t\in I:(t,t^2,\ldots,t^n)\in W_v\}
 \le\min\left\{1,\frac{n(n+1)}{v+1}\right\}.
\]
The same dimension describes the corresponding subset of the moment curve, since its parametrization is bi-Lipschitz on $I$. The reported covering argument considers each exact shell $\|q\|_\infty=Q$: there are $O(Q^{n-1})$ choices of $q$, $O(Q)$ relevant integers $p$, and each resonance sublevel
\[
 \left\{t\in I:\left|\sum_{j=1}^nq_jt^j-p\right|<Q^{-v}\right\}
\]
is covered by at most $2n+1$ intervals of diameter $O(Q^{-(v+1)/n})$. The Hausdorff--Cantelli sum is bounded by a constant times $\sum_Q Q^{n-s(v+1)/n}$, yielding the displayed coarse bound. The aid explicitly identifies this as a Remez-type elementary argument, not the requested dynamics proof or a sharp general estimate.

The text beginning ``5.3'' belongs to the next question's multiplicative-approximation setup. With $\langle x\rangle=\operatorname{dist}(x,\Z)$, it reports almost-everywhere positivity of
\[
 \liminf_{q\to\infty}q(\log q)^{2+\varepsilon}
       \langle q\alpha\rangle\langle q\beta\rangle
\]
for every $\varepsilon>0$, and attributes the stronger limit $+\infty$ to [184]. These are separate contextual measure claims, not hypotheses or conclusions of Question 26.
\subsection{Short English statement}
Estimate the Hausdorff dimension of dual $v$-approximable points on a nondegenerate manifold, preferably using dynamics. The supplied moment-curve bound is explicit but coarse and elementary; it does not establish the proposed exact formula for all manifolds.
\subsection{Sources}
\begin{itemize}
\item D. Kleinbock, Question 26, American Institute of Mathematics, \emph{Emerging applications of measure rigidity}, June 2004, section 5.2, \url{https://aimath.org/WWN/measrigid/measrigid.pdf}. The question cites [41]; the following section 5.3 cites [184]. Full bibliography entries are not supplied.
\item Companion local record, permanent ID $20003299$, source code \texttt{AIM-OTHER-0106}, canonical \texttt{aim-other-notes.json}, \texttt{notes[105]}, supplies the dual max-norm definition and chart nondegeneracy convention, referring to [109] and [17].
\item Local UnsolvedMath registry, permanent ID $20003301$, source code \texttt{catalogue entry}, canonical \texttt{aim-other-notes.json}, \texttt{notes[107]}, dated 14 August 2026, statement status reconstructed and unverified. Research attempt 1: partial result, proved rigor, low novelty confidence; machine-generated, not peer reviewed. Exact modern dimension results are reported without full citations.
\end{itemize}
\end{document}
